\section{Actual tangent and point conditions at coefficient prime \(269\)}
\label{sec:cubic-269-model-conditions}

Use the actual global ring \(S\), integral basepoint \(x_0\),
and graded model of AC-G65-model, with all of that theorem's
quantified source inputs. Put \(p=269\). The following statements
extend specific formal and local conditions to this ring. They do
not extend the mixed Euler-system vanishing at the ramified CM prime.

\begin{proposition}[AC-G66-BK: every actual dual tangent has the required local condition]
\label{AC-G66-BK}
Let \(F:S\to\overline{\mathbb Q}_p[\epsilon]/(\epsilon^2)\)
be an actual algebraic dual point, with first component \(f\).
The matrix cocycle
\[
 z_F(g)=dF(r^{\rm univ}(g).\mathrm{mat})
                 f(r^{\rm univ}(g).\mathrm{mat})^{-1}
                                                        \tag{RCB.1}
\]
is continuous over a finite coefficient extension and belongs
to the literal source \texttt{Af.IsBK} subspace for
\(\operatorname{ad}(f(r^{\rm univ}))\).
No \texttt{M2} premise is required.
\end{proposition}
\begin{proof}
The argument of AC-G62-cont applies to this actual complete
Noetherian \(O\)-algebra without using any prime-range hypothesis.
For clarity, \(f(S)\) is finite over \(O\) by the complete-local
finite-image theorem used in AC-G58. If \(J=\ker f\), the
derivation kills \(J^2\), while \(J/J^2\) is finite over
\(S/J\). Thus the dual map has coefficients in one finite
extension and is bounded on a finite \(O\)-module of generators
of \(S/J^2\). Products of sufficiently many maximal-ideal
generators tend to zero under both components, by the product
rule and the ultrametric inequality. This proves the same
continuity assertion as AC-G62-cont. The semidirect product
rule, whose differential outside \(G_K\) is
\(X\mapsto-JX^tJ\), gives exactly the cocycle identity in
(RCB.1).

At a rational prime \(v\ne p\), the actual ring kills
\(I_v^0=I_v\cap\ker\bar r_*\), and so do the representation
and its derivative. The cocycle factors on inertia through
the finite group \(I_v/I_v^0\). Its class is zero: averaging
\(z_F(h)\) over that finite group gives a cobounding vector
in the coefficient field. This step is rational and asserts
no integral denominator bound.

At \(p\), AC-G65-local applies to the entire finite
dual-number point, not just to its reduction. On
\(H=G_{K_p}\) it is potentially crystalline of the fixed
type. Pass to a finite Galois extension \(L/\mathbb Q_p\)
over which its linear restriction is crystalline, enlarging
a preliminary extension to its Galois closure if necessary.
The dual representation is then a crystalline extension of
the specialized rank-three representation by itself.
The crystalline period functor is exact on this extension;
equivalently its full period dimension supplies a
\(B_{\rm cris}\)-linear Galois-equivariant splitting.
Thus the restriction of (RCB.1) is a coboundary in
\(B_{\rm cris}\otimes\operatorname{ad}(r_f)\).
This is the same period-dimension argument proved in
AC-G47-crystalline and used in AC-G53-BK.

Restriction on continuous cohomology with this period-ring
coefficient is injective from \(G_{\mathbb Q_p}\) to
\(G_L\): restriction followed by corestriction is
\([L:\mathbb Q_p]\), an invertible scalar. One may equivalently
average the resulting finite-quotient cocycle after a
\(G_L\)-coboundary is removed. The coefficient here is the
actual linear adjoint of the semidirect representation, so
this descent includes the nontrivial quadratic coset.
It proves the required \(B_{\rm cris}\) local coboundary on
the rational decomposition group. Together with the
away-inertia calculation this is precisely \texttt{Af.IsBK}.
\end{proof}

\begin{theorem}[AC-G66-origin: basepoint rigidity and M1 at \(269\)]
\label{AC-G66-origin}
For the actual model of AC-G65, the source basepoint
condition \texttt{PtHyp.rigid0} and the full condition
\texttt{M1} hold, without \texttt{M2}.
\end{theorem}
\begin{proof}
Let a dual point at the actual basepoint have tangent of weight
\(b\leq0\). If \(b<-2\), its tangent is zero because there is
no such matrix weight. For the remaining weights, the nonpositive
adjoint pieces of the actual basepoint are
\[
\begin{array}{c|c}
-1&\bigl(\operatorname{Ind}_{G_K}^{G_{\mathbb Q}}
             E(\chi^{-j})\otimes\psi^{-1}\bigr)(-e)\\
-2&E(\eta_K\psi^{-2})(-2e)\\
0&E\oplus E(\eta_K)\oplus E(\eta_K).
\end{array}                                                \tag{RCB.2}
\]
This is the same global matrix computation as (G53.8);
ramification of the finite cubic character at the coefficient
prime does not change it. Their coordinate projections are
equivariant. AC-G66-BK puts each projected tangent in its
Bloch--Kato subspace. AC-G53-Artin applies to all these
finite-image Artin coefficients followed by nonpositive Tate
twists, with no restriction on their finite splitting field
or its degree. Hence every such tangent is a coboundary.

We also need the literal zero derivative, not only its
class modulo conjugation. The normalized cocycle vanishes
at complex conjugation. A cobounding matrix therefore has
the form
\[
 X=\begin{pmatrix}a&u&0\\v&0&-u\\0&-v&-a\end{pmatrix}.
\]
The two differentiated normalization equations at the marked
global element give \(-\Delta u=\Delta v=0\), where
\(\Delta=\chi(g_1)^j(\gamma(g_1)-\gamma(g_1)^{-1})\)
is the same unit as in AC-G53. Its unit property uses only
the retained normalized global datum; multiplication by
a cubic value is a unit also at \(269\).
Thus \(u=v=0\), and the remaining diagonal matrix is fixed
by the full adjoint action. Its coboundary is zero.
The derivative kills every universal matrix entry, hence
the dense algebra they generate, and by continuity all of
\(S\). This proves \texttt{rigid0}.

To verify the absence of an \texttt{M2} premise in the
model consequence, apply the origin-completion construction
of AC-G56 to the present model. It uses the actual complete
ring, the graded chart with the five multiplicities, and
the saturated local cotangent property. AC-G65-model and
AC-G65-local give precisely those fields, with the same
normalization equations. A derivation of weight
\(-2,-1,0\) at the origin that kills the model equations
extends there to an actual dual point, without degree-zero
finiteness. The preceding zero statement makes that
derivation zero. Duality on each finite cotangent weight
space shows that the equation classes span it rationally;
clearing their finitely many denominators gives exactly
the source's saturated \texttt{SpansCot} assertion.
Together with the already constructed smooth chart and
its positive dimension this is the full \texttt{M1}.
Every ingredient used by AC-G56 has thus been supplied for
this actual ring; no mixed tangent vanishing is used.
\end{proof}

\begin{theorem}[AC-G66-fixed: the fixed-point condition and M3 at \(269\)]
\label{AC-G66-fixed}
Every actual point of the new ring which satisfies both
\texttt{Upper} and \texttt{Lower} is the cubic basepoint.
The actual model satisfies \texttt{M3}, without \texttt{M2}.
\end{theorem}
\begin{proof}
The two triangular conditions make the representation
diagonal on \(G_K\); normalization fixes its value at
complex conjugation. The finite-image and integrality
argument of AC-G58 puts its entries in the integers of
one finite coefficient extension, with the prescribed
residual labels. Thus AC-G65-flag applies. If the residue
field has enlarged, its explicit projector proof still
applies verbatim: every denominator is the image of the
same unit in \(O\), and the coordinate reductions remain
distinct after the residue embedding.

Write its diagonal characters as \(a_i\), and set
\(b=(\chi^j\gamma,1,\chi^j\gamma^{-1})\).
At the inertia element (RC.17), the diagonal representation
has the three distinct eigenvalues \(\gamma(s),1,
\gamma(s)^{-1}\). Their residual labels fix the coordinate
assignment. Hence the ordered flag (RC.15), even though
initially expressed after strict conjugation, implies
\(a_i|_{I_p}=b_i|_{I_p}\) in the original coordinates.

At every finite place \(w\nmid p\), the actual killing
condition and the prime-to-\(p\) order of its residual
inertia image imply the same equality. More explicitly,
for \(m=3\operatorname{ord}(\psi)\), every inertia element
has its \(m\)-th power in the killed residual kernel.
Both \(a_i(g)\) and \(b_i(g)\) have \(m\)-th power one.
Their quotient reduces to one, and reduction is injective
on prime-to-\(p\) roots of unity. Thus their quotient is
one. Each \(a_i/b_i\) is consequently everywhere
unramified on \(G_K\). The class number of
\(\mathbb Q(\omega)\) is one, so global reciprocity gives
\(a_i=b_i\). Normalization at conjugation gives equality
of the full semidirect representations. Entry density and
continuity give equality of the point maps.

The all-point extension theorem of AC-G58 uses the actual
complete graded model and its saturated formal chart; its
proof does not use the unramifiedness of \(\chi\) at the
coefficient prime or \texttt{M2}. These structural fields
hold by AC-G65. Apply that extension to a model point
annihilating all positive and negative weights. Its
universal matrix entries outside the diagonal vanish, so
the resulting actual point satisfies both triangular
conditions and is the basepoint by the preceding argument.
This is the literal model fixed-origin assertion
\texttt{M3}. It is exactly the torus/entry argument of
AC-G58-M3 with the newly proved fixed-point field.
\end{proof}

\paragraph{Remaining gate.}
These results supply the actual local tangent condition,
the nonpositive origin rigidity and the fixed-point fields
at the one admissible coefficient prime left out of the
earlier good-prime statements. They do not prove rigidity
at every Type I or Type II point. In particular the mixed
coefficient at \(269\) is still outside the proved CM
Euler constructions. The global degree-zero support
condition \texttt{M2}, actual Type I Hecke-ideal containment,
and the numerical integral index estimates remain open.
The geometric objects are the same formal deformation
scheme and integral local system of AC-G65; its previously
audited SGA~1 dictionary supplies no additional vanishing
or arithmetic size theorem.
